% Wersja źródłowa: audyt II, jednostka 25
\setcounter{section}{24}
\section{Przestrzenie funkcji testowych i dystrybucji: synteza topologiczna}
\label{sec:synteza25}

\noindent\textbf{Wyniki wykorzystywane.} Korzystamy z zupełności
$\mathcal E,\mathcal S$ (Twierdzenia~\ref{thm:frechet-cinf17},
\ref{thm:schwartz-frechet17}), struktury LF i zupełności $\mathcal D$
(Twierdzenia~\ref{thm:D-stage18}, \ref{thm:D-complete18}),
własności Montela (Twierdzenie~\ref{thm:montel-test23}),
twierdzenia Banacha--Steinhausa oraz identyfikacji iloczynów
tensorowych z rozdziału~\ref{sec:tensor24}.
\par\medskip

Po omówieniu przestrzeni lokalnie wypukłych, dystrybucji i produktów
 tensorowych możemy porównać najważniejsze przestrzenie funkcyjne bez
 utożsamiania ich topologii. W tym rozdziale $d\ge1$ jest liczbą całkowitą, $\K=\C$ lub
 $\K=\R$, a $\Omega\subseteq\R^d$ jest niepustym zbiorem otwartym.
 Funkcje testowe i dystrybucje mogą mieć wartości rzeczywiste lub
 zespolone; w twierdzeniach dotyczących transformaty Fouriera
 przyjmujemy $\K=\C$. Wszędzie $\supp f$ oznacza nośnik funkcji
 lub dystrybucji. Parowanie dystrybucji z funkcją testową
 jest liniowe również nad $\C$.

\subsection{Trzy przestrzenie funkcji: definicje i zanurzenia}

\begin{definicja}[Przestrzenie $\mathcal E$, $\mathcal D$ i $\mathcal S$]
\label{def:trzy-przestrzenie25}
Dla niepustego zbioru otwartego $\Omega\subseteq\R^d$ oznaczamy
$\mathcal E(\Omega)=C^\infty(\Omega)$, wyposażoną w topologię
jednostajnej zbieżności wszystkich pochodnych na zbiorach zwartych,
wyznaczoną przez seminormy
$p_{K,m}(f)=\max_{|\alpha|\le m}\sup_{x\in K}|\partial^\alpha f(x)|$,
gdzie $K\Subset\Omega$ jest niepustym zbiorem zwartym,
a $m\in\N_0$. Niech $\mathcal D_K(\Omega)$ oznacza podprzestrzeń
funkcji gładkich o nośniku zawartym w $K$, z topologią
seminorm $p_{K,m}$. Przez $\mathcal D(\Omega)=C_c^\infty(\Omega)$
rozumiemy ścisłą lokalnie wypukłą granicę induktywną
etapów $\mathcal D_{K_j}(\Omega)$, gdzie
$K_j\subseteq\operatorname{int}K_{j+1}$ i $\bigcup_jK_j=\Omega$.
Topologia nie zależy od wyczerpania
(Stwierdzenie~\ref{prop:D-independence18}).
Przestrzeń Schwartza $\mathcal S(\R^d)$ składa się z
tych $f\in C^\infty(\R^d)$, dla których
\[
 p_{N,m}(f)=\max_{|\alpha|\le m}\sup_{x\in\R^d}
 (1+|x|)^N|\partial^\alpha f(x)|<\infty
 \quad(N,m\in\N_0).
\]
Te normy wyznaczają jej topologię i odpowiadają normom
$Q_{N,m}$ z Definicji~\ref{def:schwartz17}.
Wszystkie przestrzenie funkcji rozpatrujemy nad ustalonym $\K$.
Przestrzenie dualne topologiczne oznaczamy odpowiednio
$\mathcal E'(\Omega)$, $\mathcal D'(\Omega)$ i
$\mathcal S'(\R^d)$. Jeżeli podajemy topologię na przestrzeni
 dualnej $X'$, indeks $\sigma$ oznacza topologię słabą-*
 $\sigma(X',X)$, wyznaczoną przez seminormy $T\mapsto|T(f)|$
 dla $f\in X$, zaś $\beta$ oznacza silną topologię
 $\beta(X',X)$, wyznaczoną przez seminormy
 $T\mapsto\sup_{f\in B}|T(f)|$ dla zbiorów ograniczonych $B\subseteq X$.
\end{definicja}

\begin{twierdzenie}[Ciągłe i gęste zanurzenia]
\label{thm:zanurzenia25}
Niech $d\ge1$. Odwzorowania identycznościowe
\[
 \mathcal D(\R^d)\longrightarrow\mathcal S(\R^d)
 \longrightarrow\mathcal E(\R^d)
\]
są ciągłymi, różnowartościowymi odwzorowaniami liniowymi o gęstych
obrazach. Dla dowolnego niepustego zbioru otwartego $\Omega\subseteq\R^d$
takie samo stwierdzenie zachodzi dla zanurzenia
$\mathcal D(\Omega)\to\mathcal E(\Omega)$.
Żadne z trzech wskazanych odwzorowań nie jest
izomorfizmem liniowo-topologicznym na swój obraz z topologią
podprzestrzeni; ostatnie stwierdzenie zachodzi dla każdego
niepustego zbioru otwartego $\Omega$.
\end{twierdzenie}
\begin{proof}
Dla ustalonego zwartego $K\Subset\R^d$ oraz $N,m\in\N_0$
mamy $p_{N,m}(f)\le (1+\sup_{x\in K}|x|)^N p_{K,m}(f)$
dla $f\in\mathcal D_K(\R^d)$. Stąd ciągłość
$\mathcal D_K\to\mathcal S$; własność uniwersalna granicy
induktywnej daje ciągłość $\mathcal D\to\mathcal S$.
Ciągłość $\mathcal S\to\mathcal E$ wynika z oszacowania
$p_{K,m}(f)\le p_{0,m}(f)$. Analogicznie sprawdzamy
ciągłość $\mathcal D(\Omega)\to\mathcal E(\Omega)$.

Aby wykazać gęstość, wybierzmy $\chi\in\mathcal D(\R^d)$
równe $1$ na kuli jednostkowej i zdefiniujmy
$\chi_R(x)=\chi(x/R)$ dla $R\ge1$. Dla $f\in\mathcal S(\R^d)$ i $|\alpha|\le m$
wzór Leibniza daje
\[
 \partial^\alpha((1-\chi_R)f)
 =\sum_{\gamma\le\alpha}\binom{\alpha}{\gamma}
 \partial^\gamma(1-\chi_R)\partial^{\alpha-\gamma}f.
\]
Wszystkie czynniki $\partial^\gamma(1-\chi_R)$
znikają dla $|x|<R$. Dla $\gamma\ne0$ są ograniczone
przez $R^{-|\gamma|}\sup|\partial^\gamma\chi|$,
a dla $\gamma=0$ przez $1+\sup|\chi|$.
Ponieważ $R\ge1$ i $(1+|x|)^{-1}\le(1+R)^{-1}$
na pozostałej części przestrzeni, otrzymujemy
\[
 p_{N,m}((1-\chi_R)f)
 \le \frac{C_{\chi,m}}{1+R}p_{N+1,m}(f).
\]
Zatem $\chi_Rf\to f$ w $\mathcal S$ i $\chi_Rf\in\mathcal D$.
Dla $f\in\mathcal E(\R^d)$ mamy $\chi_Rf=f$ w otoczeniu
każdego ustalonego zbioru zwartego przy dostatecznie dużym $R$,
a więc $\chi_Rf\to f$ w $\mathcal E$. W ogólnej
$\Omega$ używamy funkcji odcięcia $\chi_j\in\mathcal D(\Omega)$,
równej $1$ w otoczeniu $K_j$, gdzie $(K_j)$ jest
wyczerpaniem $\Omega$ zbiorami zwartymi. Każdy zwarty $L\Subset\Omega$ mieści się w pewnym $K_J$,
więc dla $j\ge J$ wszystkie pochodne różnicy znikają na $L$.
Stąd $\chi_j f\to f$ w $\mathcal E(\Omega)$.

Na $\R^d$ niech $\psi\in\mathcal D(\R^d)$ będzie niezerowa
i $f_n(x)=\psi(x-ne_1)$. Wówczas $f_n\to0$ w
$\mathcal E(\R^d)$, bo nośniki ostatecznie omijają każdy
ustalony zbiór zwarty, ale nie w $\mathcal S(\R^d)$,
gdyż $p_{0,0}(f_n)=p_{0,0}(\psi)>0$.
Z kolei $g_n(x)=e^{-n^2}\psi(x-ne_1)$ zbiega do zera w
$\mathcal S(\R^d)$: dla $\supp\psi\subseteq\overline B(0,A)$
mamy
\[
 p_{N,m}(g_n)\le e^{-n^2}(1+n+A)^N p_{0,m}(\psi)
 \longrightarrow0.
\]
Nie zbiega do zera w $\mathcal D(\R^d)$, gdyż nośniki
nie mieszczą się w jednym zbiorze zwartym
(Twierdzenie~\ref{thm:D-sequences18}).

Dla dowolnego niepustego otwartego $\Omega\subseteq\R^d$
wybieramy wyczerpanie $(K_j)$ jak w definicji.
Zbiór $\Omega\setminus K_j$ jest niepusty i otwarty,
bo żaden niepusty otwarty podzbiór $\R^d$ nie jest zwarty.
Wybieramy niezerowe $h_j\in\mathcal D(\Omega)$ takie,
że $\supp h_j\subseteq\Omega\setminus K_j$.
Wówczas $h_j\to0$ w $\mathcal E(\Omega)$, gdyż dla
każdego zwartego $L\Subset\Omega$ nośniki ostatecznie
omijają jego otoczenie (wybieramy $L\subseteq\operatorname{int}K_J$).
Gdyby $h_j\to0$ w $\mathcal D(\Omega)$, wszystkie nośniki
byłyby zawarte w pewnym zwartym $L\Subset\Omega$,
a więc w pewnym $K_J$. Dla $j\ge J$ oznaczałoby to
$h_j=0$, sprzecznie z wyborem. Topologia $\mathcal D(\Omega)$
jest zatem ściśle silniejsza od odziedziczonej z $\mathcal E(\Omega)$.
\end{proof}

\begin{uwaga}[Zupełność i lokalna wypukłość]
\label{rem:tabela-struktur25}
$\mathcal E(\Omega)$ i $\mathcal S(\R^d)$ są nuklearnymi
przestrzeniami Frécheta, a $\mathcal D(\Omega)$ jest nuklearną,
zupełną, niemetryzowalną przestrzenią LF
(Twierdzenia~\ref{thm:D-complete18}, \ref{thm:D-nonmetric18}).
Wszystkie trzy są przestrzeniami Montela
(Twierdzenie~\ref{thm:montel-test23}).
Nuklearność $\mathcal E$ i $\mathcal S$ wynika odpowiednio
z Twierdzenia~\ref{thm:nuclear-cinfty23}
i Wniosku~\ref{cor:nuclear-schwartz23}.
Trwałość nuklearności względem podprzestrzeni, ilorazów
i przeliczalnych granic induktywnych udowodniono w
Twierdzeniach~\ref{thm:nuclear-subquot23} i~\ref{thm:nuclear-indlim23};
nuklearność $\mathcal D_K$ i $\mathcal D$ wynika z
Twierdzenia~\ref{thm:nuclear-D23}. Nie należy mylić nuklearności
z własnością Montela: te dwie własności były w skrypcie
uzasadniane oddzielnie.
\end{uwaga}

\subsection{Przechodzenie do przestrzeni dualnych topologicznych}

\begin{lemat}[Transpozycja gęstego zanurzenia]
\label{lem:transpozycja25}
Niech $X,Y$ będą przestrzeniami lokalnie wypukłymi
Hausdorffa, a $j:X\to Y$ ciągłym odwzorowaniem liniowym
o gęstym obrazie. Wówczas transpozycja
$j':Y'\to X'$, $j'(T)=T\circ j$, jest liniowa,
różnowartościowa i ciągła zarówno dla topologii słabych-*, jak i dla
silnych topologii przestrzeni dualnych.
\end{lemat}
\begin{proof}
Dla każdego $T\in Y'$ złożenie $T\circ j$ jest ciągłym
funkcjonałem liniowym na $X$, więc $j'$ jest dobrze
określone; liniowość wynika z liniowości składania
względem pierwszego czynnika. Jeżeli $T\circ j=0$,
to $T=0$ przez gęstość $j(X)$ i ciągłość $T$.
Dla $x\in X$ mamy
$|\langle j'T,x\rangle|=|\langle T,jx\rangle|$,
co daje ciągłość w topologiach słabych-*.
Jeśli $B\subseteq X$ jest ograniczony, to $j(B)$ jest
ograniczony w $Y$ (dla każdej ciągłej seminormy $p$
na $Y$ seminorma $p\circ j$ jest ciągła na $X$).
Dla seminorm silnych zachodzi dokładna równość
$\sup_{x\in B}|\langle j'T,x\rangle|
= \sup_{y\in j(B)}|\langle T,y\rangle|$.
To dowodzi ciągłości w topologiach silnych.
\end{proof}

\begin{wniosek}[Łańcuch przestrzeni dualnych]
\label{cor:dual-lancuch25}
Dla każdego $d\ge1$ zachodzą kanoniczne, ciągłe i
różnowartościowe odwzorowania liniowe
\[
 \mathcal E'(\R^d)\longrightarrow\mathcal S'(\R^d)
 \longrightarrow\mathcal D'(\R^d),
\]
zarówno w topologiach słabych-*, jak i silnych.
Dla każdego niepustego zbioru otwartego $\Omega\subseteq\R^d$
kanoniczne odwzorowanie
$\mathcal E'(\Omega)\to\mathcal D'(\Omega)$
ma te same własności. Samo stwierdzenie ciągłości i
różnowartościowości \emph{nie oznacza}, że są to zanurzenia
topologiczne na obrazy z topologią podprzestrzeni.
\end{wniosek}
\begin{proof}
Stosujemy Lemat~\ref{lem:transpozycja25} do odwzorowań
z Twierdzenia~\ref{thm:zanurzenia25}. Jeżeli $i:X\to Y$
oraz $j:Y\to Z$ są tymi kolejnymi odwzorowaniami,
to $(j\circ i)'=i'\circ j'$, gdyż dla $T\in Z'$
mamy $(T\circ j)\circ i=T\circ(j\circ i)$.
Dlatego ograniczenie z $\mathcal E'$ do $\mathcal D'$
przez $\mathcal S'$ jest tym samym odwzorowaniem co
bezpośrednie ograniczenie. W przypadku zespolonym
transpozycja jest liniowa, bez sprzężenia zespolonego.
Przykłady nierównoważności topologii na obrazach
podamy w Przykładzie~\ref{ex:escaping-deltas25}.
\end{proof}

\begin{twierdzenie}[Dystrybucje o zwartym nośniku]
\label{thm:compact-support-E25}
Niech $\Omega\subseteq\R^d$ będzie niepustym zbiorem otwartym.
Obraz kanonicznego odwzorowania
$\mathcal E'(\Omega)\to\mathcal D'(\Omega)$
składa się dokładnie z dystrybucji o zwartym nośniku.
Dokładniej, dla każdego $F\in\mathcal E'(\Omega)$ istnieją
niepusty zwarty $K\Subset\Omega$, $m\in\N_0$ i $C>0$ takie, że
\[
 |F(f)|\le C\max_{|\alpha|\le m}\sup_{x\in K}
 |\partial^\alpha f(x)|\qquad(f\in\mathcal E(\Omega)).
\]
Odwrotnie, jeżeli $T\in\mathcal D'(\Omega)$ ma zwarty nośnik,
to istnieje dokładnie jeden funkcjonał ciągły
$F\in\mathcal E'(\Omega)$ równy $T$ na $\mathcal D(\Omega)$.
\end{twierdzenie}
\begin{proof}
Z kryterium ciągłości funkcjonału na przestrzeni
$\mathcal E(\Omega)$ (Twierdzenie~\ref{thm:dual-cinf17})
otrzymujemy oszacowanie przez skończenie wiele
seminorm $p_{K_j,m_j}$. Biorąc $K=\bigcup_jK_j$
i $m=\max_jm_j$, otrzymujemy oszacowanie z tezy
(dla funkcjonału zerowego można wybrać dowolny
niepusty zwarty $K$). Ograniczenie
$T=F|_{\mathcal D(\Omega)}$ jest dystrybucją na mocy
Wniosku~\ref{cor:dual-lancuch25}. Gdy
$\varphi\in\mathcal D(\Omega\setminus K)$, jej
przedłużenie przez zero wraz ze wszystkimi pochodnymi
znika w otoczeniu $K$. Zatem $T(\varphi)=0$,
czyli $\supp T\subseteq K$.
Nośnik jest tu nośnikiem dystrybucji $T$,
a nie zbiorem punktów przestrzeni argumentów funkcjonału $F$.

Niech teraz $K=\supp T\Subset\Omega$. Jeśli $K$ jest
pusty, to $T=0$ z lokalności dystrybucji
(Lemat~\ref{lem:distribution-local19}) i wystarcza $F=0$.
W przeciwnym przypadku wybierzmy
$\chi\in\mathcal D(\Omega)$ równą $1$ w otoczeniu $K$
(Lemat~\ref{lem:cutoffs17}) i połóżmy $L=\supp\chi$.
Kładziemy $F(f)=\langle T,\chi f\rangle$ dla
$f\in\mathcal E(\Omega)$. Ponieważ $\supp(\chi f)\subseteq L$,
z lokalnego oszacowania dystrybucji
(Twierdzenie~\ref{thm:distribution-estimate19})
istnieją $C_L>0$ i $m\in\N_0$ takie, że
\[
 |F(f)|\le C_L p_{L,m}(\chi f)
 \le C_L 2^m p_{L,m}(\chi)p_{L,m}(f).
\]
Ostatnia nierówność wynika z reguły Leibniza i tożsamości
$\sum_{\gamma\le\alpha}\binom{\alpha}{\gamma}=2^{|\alpha|}$.
Stąd $F$ jest liniowym funkcjonałem ciągłym.
Wartość nie zależy od wyboru $\chi$: dla dwóch odcięć
$\chi_1,\chi_2$ funkcja $(\chi_1-\chi_2)f$ jest testowa
i znika w otoczeniu $K$. Jej nośnik jest więc zwartym
podzbiorem $\Omega\setminus K$. Z definicji nośnika
i lokalności dystrybucji wynika, że $T$ ją anihiluje. Jeżeli $f\in\mathcal D(\Omega)$, to
$T(f)=T(\chi f)$, ponieważ $(1-\chi)f$ znika w
otoczeniu $K$. Jednoznaczność przedłużenia wynika
z gęstości $\mathcal D(\Omega)$ w $\mathcal E(\Omega)$.
Otrzymana identyfikacja jest zgodna z przedłużeniem
z Lematu~\ref{lem:compact-action21} oraz Uwagą~
\ref{rem:Eprime21}; nie nadajemy obrazowi automatycznie
topologii podprzestrzeni $\mathcal D'$.
\end{proof}

\subsection{Dystrybucje temperowane jako podprzestrzeń \texorpdfstring{$\mathcal D'$}{D'}}

\begin{twierdzenie}[Kryterium temperowania]
\label{thm:temp-criterion25}
Niech $d\ge1$ będzie liczbą całkowitą, $\K\in\{\R,\C\}$
i $T\in\mathcal D'(\R^d)$ nad $\K$.
Następujące warunki są równoważne:
\begin{enumerate}[label=\textup{(\roman*)}]
\item $T$ należy do obrazu kanonicznego odwzorowania
$\mathcal S'(\R^d)\to\mathcal D'(\R^d)$;
\item istnieją $C>0$, $N,m\in\N_0$ takie, że
$|\langle T,\varphi\rangle|\le C p_{N,m}(\varphi)$
dla każdej $\varphi\in\mathcal D(\R^d)$.
\end{enumerate}
Jeżeli te warunki zachodzą, przedłużenie
$\widetilde T\in\mathcal S'(\R^d)$ jest jednoznaczne
i spełnia $|\langle\widetilde T,f\rangle|\le C p_{N,m}(f)$
dla wszystkich $f\in\mathcal S(\R^d)$, z tymi samymi
stałymi co w~\textup{(ii)}. W szczególności każda
dystrybucja o zwartym nośniku jest temperowana.
\end{twierdzenie}
\begin{proof}
Niech najpierw $\widetilde T\in\mathcal S'$ będzie
przedłużeniem $T$. Z ciągłości w zerze i współkońcowości
rosnącego ciągu norm $p_{j,j}$ istnieją $j\in\N_0$ oraz
$\delta>0$ takie, że $p_{j,j}(f)<\delta$ implikuje
$|\widetilde T(f)|<1$. Dla $f\ne0$ stosujemy ten warunek
do $\delta f/(2p_{j,j}(f))$, otrzymując
$|\widetilde T(f)|\le (2/\delta)p_{j,j}(f)$.
Dla $f=0$ nierówność jest oczywista. Ograniczenie
do $\mathcal D$ daje warunek~\textup{(ii)}.
Jest to argument z Twierdzenia~\ref{thm:tempered-estimate22},
działający nad obydwoma ciałami.

Załóżmy teraz~\textup{(ii)}. Dla każdego $f\in\mathcal S$
wybieramy ciąg $(\varphi_j)$ w $\mathcal D$ zbieżny do $f$
w $\mathcal S$ (Twierdzenie~\ref{thm:zanurzenia25}).
Ponieważ
\[
 |T(\varphi_j)-T(\varphi_k)|
 \le C p_{N,m}(\varphi_j-\varphi_k)\longrightarrow0
 \quad(j,k\to\infty),
\]
ciąg wartości ma granicę w zupełnym ciele $\K$.
Kładziemy $\widetilde T(f)=\lim_jT(\varphi_j)$.
Dla drugiego ciągu $(\psi_j)$ zbieżnego do $f$ mamy
$|T(\varphi_j)-T(\psi_j)|\le
C p_{N,m}(\varphi_j-\psi_j)\to0$, więc definicja nie zależy
od wyboru ciągu. Dla $f,g\in\mathcal S$ i $a,b\in\K$
można aproksymować $af+bg$ przez $a\varphi_j+b\psi_j$,
gdzie teraz $\varphi_j\to f$ i $\psi_j\to g$.
Przechodząc do granicy, otrzymujemy liniowość.
Ponadto $|p_{N,m}(\varphi_j)-p_{N,m}(f)|
\le p_{N,m}(\varphi_j-f)$, więc
\[
 |\widetilde T(f)|\le C p_{N,m}(f).
\]
Daje to ciągłość; na $\mathcal D$ wystarczy wziąć
ciąg stały, aby uzyskać $\widetilde T|_{\mathcal D}=T$.
Jednoznaczność wynika z gęstości $\mathcal D$.

Jeżeli $T$ ma zwarty nośnik, to z
Twierdzenia~\ref{thm:compact-support-E25} otrzymujemy
$|T(\varphi)|\le C p_{K,m}(\varphi)\le C p_{0,m}(\varphi)$
na funkcjach testowych, a zatem~\textup{(ii)} zachodzi z $N=0$.
Przedłużenie jest ograniczeniem funkcjonału na $\mathcal E$
do $\mathcal S$, przez jednoznaczność; odpowiada to
Stwierdzeniu~\ref{prop:schwartz-embeddings22}.
\end{proof}

\begin{przyklad}[Dystrybucja nietemperowana]
\label{ex:non-tempered25}
Na $\R$ funkcja $f(x)=e^x$ jest lokalnie całkowalna,
więc wyznacza dystrybucję regularną $T_f\in\mathcal D'(\R)$.
Nie jest ona temperowana. Wybierzmy nieujemną, niezerową
$\psi\in\mathcal D((-1,1))$, przedłużoną zerem na $\R$,
i niech $\varphi_n(x)=e^{-n/2}\psi(x-n)$ dla $n\ge1$.
Dla $x\in\supp\varphi_n$ mamy $|x|\le n+1$; pochodne
$\varphi_n$ znikają poza tym nośnikiem. Stąd dla wszystkich
$N,m\in\N_0$
\[
 p_{N,m}(\varphi_n)\le e^{-n/2}(n+2)^N p_{0,m}(\psi)
 \longrightarrow0,
\]
gdyż funkcja wykładnicza rośnie szybciej niż każdy
wielomian. Jednak po podstawieniu $y=x-n$ otrzymujemy
\[
 \langle T_f,\varphi_n\rangle=
 e^{n/2}\int_\R e^y\psi(y)\,\mathrm d\lambda_1(y)
 \longrightarrow+\infty.
\]
Całka jest skończona i ściśle dodatnia: $\psi$ jest
nieujemna i dodatnia na pewnym niepustym przedziale.
Ciągłość przedłużenia na $\mathcal S$ wymagałaby
$\langle T_f,\varphi_n\rangle\to0$, co daje sprzeczność.
Przykład działa nad $\R$ i $\C$; razem z
Przykładem~\ref{ex:non-tempered22} pokazuje on dwie
metody wykrywania braku temperowalności (przesuwanie
nośnika oraz powiększanie odcięcia).
\end{przyklad}

\subsection{Topologie, zbieżność i charakterystyczne przykłady}

\begin{twierdzenie}[Zbieżność ciągów w przestrzeniach dualnych]
\label{thm:weak-strong-sequences25}
Niech $X$ będzie beczkową przestrzenią lokalnie wypukłą
Hausdorffa nad $\K$, w której każdy zbiór ograniczony
ma zwarte domknięcie (w szczególności przestrzenią Montela).
Dla ciągu $(T_n)$ w $X'$ i $T\in X'$ następujące
warunki są równoważne: $T_n\to T$ w topologii
$\sigma(X',X)$ oraz $T_n\to T$ w topologii $\beta(X',X)$.
Stwierdzenie stosuje się w szczególności do
$X=\mathcal D(\Omega)$, $X=\mathcal E(\Omega)$ i
$X=\mathcal S(\R^d)$.
\end{twierdzenie}
\begin{proof}
Zbieżność silna implikuje słabą-*, gdyż zbiory
jednoelementowe są ograniczone. W przeciwną stronę,
zbieżność skalarna $\langle T_n,x\rangle\to\langle T,x\rangle$
dla każdego $x$ daje punktową ograniczoność rodziny
$\{T_n-T:n\ge1\}$. Z beczkowości i twierdzenia
Banacha--Steinhausa wynika jej równociągłość.
Ustalmy ograniczony zbiór $B\subseteq X$ i $\varepsilon>0$.
Dla $B=\varnothing$ odpowiednia seminorma jest zerowa.
W przeciwnym przypadku wybieramy najpierw otwarte
otoczenie zera $U$ takie, że
$|\langle T_n-T,u\rangle|<\varepsilon/3$
dla wszystkich $n$ i $u\in U$.
Zwartość $K=\overline B^{\,X}$ pozwala wybrać skończenie
wiele punktów $x_1,\ldots,x_k\in K$ takich, że
$K\subseteq\bigcup_{j=1}^k(x_j+U)$.
Ze zbieżności punktowej istnieje $n_0$ takie, że
$|\langle T_n-T,x_j\rangle|<\varepsilon/3$
dla wszystkich $n\ge n_0$ i $j=1,\ldots,k$.
Dla $x\in B$ wybieramy $j$ i $u\in U$ z $x=x_j+u$.
Nierówność trójkąta daje
$|\langle T_n-T,x\rangle|<2\varepsilon/3$.
W konsekwencji
$\sup_{x\in B}|\langle T_n-T,x\rangle|\le2\varepsilon/3<\varepsilon$
dla $n\ge n_0$.
Zatem zbieżność jest silna. Zastosowania do $\mathcal D$,
$\mathcal E$ i $\mathcal S$ wynikają z
Twierdzenia~\ref{thm:montel-test23}.
Jest to wspólna postać Twierdzenia~
\ref{thm:weak-strong-sequences19} i Wniosku~
\ref{cor:montel-dual-seq23}.
\end{proof}

\begin{uwaga}[Zbieżność ciągów a równość topologii]
\label{rem:sequence-net25}
Twierdzenie~\ref{thm:weak-strong-sequences25} nie oznacza
równości topologii $\sigma$ i $\beta$.
W dowodzie używamy punktowej ograniczoności całej
rodziny wyrazów ciągu zbieżnego. Dla sieci zbieżnej
punktowo nie wynika stąd ograniczoność wszystkich
jej wyrazów.

Dla każdej z trzech rozważanych przestrzeni $X$
istnieje ograniczony zbiór $B_0$ zawierający nieskończenie
wiele liniowo niezależnych funkcji. Dla $\mathcal D$
jest to zbiór skonstruowany w dowodzie Twierdzenia~
\ref{thm:weak-strong-different19}; dla $\mathcal E$
i $\mathcal S$ korzystamy z jego obrazu przez ciągłe
włożenie z Twierdzenia~\ref{thm:zanurzenia25}.
Niech $I$ będzie zbiorem par $(F,n)$, gdzie
$F\subset X$ jest skończony, $n\in\N$, uporządkowanym
przez inkluzję pierwszego czynnika i zwykły porządek
drugiego. Dla każdej pary wybieramy
$u\in B_0\setminus\operatorname{lin}F$.
Skończenie wymiarowa podprzestrzeń $\operatorname{lin}F$
jest domknięta, więc twierdzenie Hahna--Banacha
zapewnia $T_{F,n}\in X'$ z
$T_{F,n}|_{\operatorname{lin}F}=0$ i $T_{F,n}(u)=n$.
Dla ustalonego $x\in X$ od indeksu $(\{x\},1)$
wszystkie wartości $T_{F,n}(x)$ są zerowe.
Sieć zbiega więc słabo-* do zera. Jednak
$\sup_{v\in B_0}|T_{F,n}(v)|\ge n$, więc
nie zbiega silnie do zera. Topologia silna
jest zatem ściśle silniejsza w każdym z trzech przypadków.
\end{uwaga}

\begin{przyklad}[Zbieżność i ucieczka nośnika]
\label{ex:escaping-deltas25}
Na $\R$ niech $\delta_n$ oznacza deltę Diraca w punkcie
$n$ dla $n\ge1$. Każda z tych dystrybucji ma zwarty
nośnik, więc należy także do $\mathcal E'$ i $\mathcal S'$
na mocy Twierdzeń~\ref{thm:compact-support-E25}
i~\ref{thm:temp-criterion25}.
Wtedy $\delta_n\to0$ silnie w $\mathcal D'(\R)$:
dla ograniczonego $B\subseteq\mathcal D(\R)$ nośniki
wszystkich jego elementów mieszczą się w jednym zwartym
przedziale (Twierdzenie~\ref{thm:D-bounded18}), więc
$\sup_{\varphi\in B}|\varphi(n)|=0$ dla dużych $n$.
Również $\delta_n\to0$ silnie w $\mathcal S'(\R)$, bo dla
ograniczonego zbioru $B\subseteq\mathcal S$ mamy
$C_B=\sup_{\varphi\in B}p_{2,0}(\varphi)<\infty$ oraz
$\sup_{\varphi\in B}|\varphi(n)|\le C_B(1+n)^{-2}$.
Jednak $\delta_n\not\to0$ nawet słabo-* w
$\mathcal E'(\R)$, ponieważ $\langle\delta_n,1\rangle=1$.

Z kolei $T_n=e^n\delta_n$ zbiega silnie do zera w
$\mathcal D'(\R)$ (na każdym ograniczonym zbiorze
seminorma silna jest od pewnego indeksu zerowa).
Nie zbiega do zera nawet słabo-* w $\mathcal S'(\R)$.
Niech $r(x)=\sqrt{1+x^2}$ i $\varphi(x)=e^{-r(x)/2}$.
Wykażemy, że $\varphi\in\mathcal S(\R)$.
Dla każdego $k\in\N_0$ istnieje wielomian $P_k(x,z)$
taki, że
\[
 \varphi^{(k)}(x)=e^{-r(x)/2}r(x)^{-2k}P_k(x,r(x)).
\]
Dla $k=0$ przyjmujemy $P_0=1$. Ponieważ $r'=x/r$,
różniczkowanie wzoru daje rekurencję
\[
 P_{k+1}=z^2\partial_x P_k+xz\partial_z P_k
          -2kxP_k-\tfrac12 xzP_k.
\]
Zatem wszystkie $P_k$ są wielomianami. Ponieważ
$1\le r(x)\le1+|x|$ i $r(x)\ge|x|$, każda pochodna
jest ograniczona przez stałą razy wielomianowy wzrost
$(1+|x|)^{M_k}$ pomnożony przez $e^{-|x|/2}$.
Wszystkie seminormy Schwartza są więc skończone.
Jednak dla $n\ge1$ mamy $r(n)\le n+1$, toteż
$\langle T_n,\varphi\rangle=e^{n-r(n)/2}
\ge e^{(n-1)/2}\to+\infty$.

Dla dowolnego niepustego zbioru otwartego $\Omega$
wybierzmy wyczerpanie $(K_j)$ i punkty
$a_j\in\Omega\setminus K_j$.
Wówczas $\delta_{a_j}\to0$ silnie w $\mathcal D'(\Omega)$,
gdyż każdy ograniczony zbiór testów ma wspólny zwarty
nośnik zawarty w pewnym $K_J$.
Natomiast $\langle\delta_{a_j},1\rangle=1$ na
$\mathcal E(\Omega)$ dla wszystkich $j$.
Zatem wszystkie odwzorowania z
Wniosku~\ref{cor:dual-lancuch25} nie są izomorfizmami
topologicznymi na swoje obrazy z topologią podprzestrzeni,
ani dla topologii słabych-*, ani dla silnych.
\end{przyklad}

\subsection{Operatory naturalne i granice ich działania}

\begin{stwierdzenie}[Różniczkowanie i mnożenie]
\label{prop:ops25}
Niech $\Omega\subseteq\R^d$ będzie niepustym zbiorem otwartym.
Dla każdego wielowskaźnika $\alpha\in\N_0^d$
operator $\partial^\alpha$ działa liniowo i ciągle na
$\mathcal E(\Omega)$ i $\mathcal D(\Omega)$, a przez
transpozycję także na $\mathcal E'(\Omega)$ i
$\mathcal D'(\Omega)$ w topologiach słabych-* i silnych.
Dla $a\in C^\infty(\Omega)$ mnożenie $f\mapsto af$
i odpowiadające mu mnożenie dystrybucji działają ciągle
na tych przestrzeniach.
Na $\mathcal S(\R^d)$ i $\mathcal S'(\R^d)$ te same
własności mnożenia zachodzą, jeżeli dla każdego
wielowskaźnika $\alpha$ istnieją $C_\alpha>0$ i
$N_\alpha\in\N_0$ takie, że
$|\partial^\alpha a(x)|\le C_\alpha(1+|x|)^{N_\alpha}$
dla wszystkich $x\in\R^d$.
\end{stwierdzenie}
\begin{proof}
Ciągłość operatorów na przestrzeniach testowych wynika z
nierówności Leibniza i z definicji seminorm. W przestrzeni
$\mathcal D$ mnożenie nie powiększa nośnika, więc
wystarcza sprawdzenie ciągłości na każdym etapie
$\mathcal D_K$. Dla $\mathcal S$ ustalmy $N,m$.
Po zastosowaniu wzoru Leibniza każdą seminormę
$p_{N,m}(af)$ można oszacować przez
$C_{N,m,a}p_{N+M,m}(f)$, gdzie $M$ jest maksimum
wykładników wzrostu pochodnych $a$ rzędu co najwyżej $m$.
Operator transponowany działa ciągle w topologii słabej-*,
a w silnej dlatego, że ciągłe odwzorowania liniowe
przenoszą zbiory ograniczone na ograniczone. Dla
pochodnej dystrybucyjnej pojawia się zwykły czynnik
$(-1)^{|\alpha|}$ w formule transpozycji.
\end{proof}

\begin{przyklad}[Nie każdy mnożnik gładki zachowuje przestrzeń Schwartza]
\label{ex:not-multiplier25}
Funkcja $a(x)=e^{|x|^2}$ jest gładka, ale mnożenie przez nią
nie zachowuje przestrzeni Schwartza: dla
$f(x)=e^{-|x|^2}\in\mathcal S(\R^d)$ iloczyn $af=1$
nie należy do $\mathcal S(\R^d)$. W przestrzeniach
$\mathcal D$ i $\mathcal E$ mnożenie przez $a$ nie
powoduje takiego problemu.
\end{przyklad}

\begin{uwaga}[Transformacja Fouriera i splot]
\label{rem:fourier-conv25}
Transformacja Fouriera jest automorfizmem
$\mathcal S(\R^d)$ i $\mathcal S'(\R^d)$.
Nie jest automorfizmem $\mathcal D'(\R^d)$, ponieważ
nie każda dystrybucja należy do $\mathcal S'$, a jej
transformata Fouriera nie musi być zdefiniowana przez
zwykłą transpozycję na przestrzeni testowej $\mathcal D$.
Splot z dystrybucją o zwartym nośniku jest natomiast
określony zarówno na $\mathcal D'$ jak i na $\mathcal S'$;
splot dwóch dowolnych dystrybucji nie musi być określony.
Teoria mnożenia dystrybucji przez funkcje gładkie i splotu
została udowodniona we wcześniejszych rozdziałach.
\end{uwaga}

\subsection{Lokalność i problemy przedłużania dystrybucji}

\begin{twierdzenie}[Ograniczenie i brak automatycznego przedłużenia]
\label{thm:restriction25}
Jeżeli $U\subseteq\Omega\subseteq\R^d$ są niepustymi
zbiorami otwartymi, to odwzorowanie ograniczenia
$\mathcal D'(\Omega)\to\mathcal D'(U)$ jest liniowe
i ciągłe w topologiach słabych-* oraz silnych.
Nie każda dystrybucja na $U$ przedłuża się do
dystrybucji na $\Omega$.
\end{twierdzenie}
\begin{proof}
Funkcję testową $\varphi\in\mathcal D(U)$ przedłużamy
przez zero do $\Omega$ i oznaczamy przedłużenie przez
$\widetilde\varphi$. Operator
$\mathcal D(U)\to\mathcal D(\Omega)$,
$\varphi\mapsto\widetilde\varphi$, jest ciągły: na każdym
etapie z ustalonym zwartym nośnikiem wynika to z równości
pochodnych i seminorm, a na całej przestrzeni z własności
uniwersalnej granicy induktywnej. Stąd transpozycja
daje ciągłe ograniczenie w obydwu topologiach dualnych.

Aby wykazać drugie stwierdzenie, weźmy
$U=(0,1)$, $\Omega=(-1,1)$ i dystrybucję regularną
$T_f$ wyznaczoną przez $f(x)=\exp(1/x^2)$ na $U$.
Gdyby miała przedłużenie $\widetilde T\in\mathcal D'(\Omega)$,
to dla pewnego zwartego $K=[-1/2,1/2]\Subset\Omega$,
$m\in\N_0$ i $C>0$ zachodziłoby
$|\langle\widetilde T,\varphi\rangle|\le C p_{K,m}(\varphi)$
dla wszystkich $\varphi\in\mathcal D_K(\Omega)$.
Niech $\psi\in\mathcal D((-1,1))$ będzie nieujemna i niezerowa,
$r_n=1/(10n^2)$ oraz
$\varphi_n(x)=\psi((x-1/n)/r_n)$ dla dużych $n$.
Wówczas $\supp\varphi_n\subset U\cap K$,
$p_{K,m}(\varphi_n)=O(n^{2m})$, natomiast
$\langle T_f,\varphi_n\rangle\ge c n^{-2}e^{c'n^2}$
dla pewnych $c,c'>0$. To sprzeczność.
\end{proof}

\subsection{Schemat zależności i produkty tensorowe}

\begin{stwierdzenie}[Zestawienie kanonicznych odwzorowań]
\label{prop:summary25}
Dla każdego $d\ge1$ istnieją ciągłe i gęste zanurzenia
liniowe $\mathcal D(\R^d)\to\mathcal S(\R^d)\to
\mathcal E(\R^d)$ oraz dualne odwzorowania ciągłe
i różnowartościowe
$\mathcal E'(\R^d)\to\mathcal S'(\R^d)\to
\mathcal D'(\R^d)$. Dla otwartych niepustych
$\Omega\subseteq\R^{d_1}$, $\Lambda\subseteq\R^{d_2}$
kanoniczne izomorfizmy
\[
\mathcal E(\Omega)\widehat\otimes_\pi\mathcal E(\Lambda)
\simeq\mathcal E(\Omega\times\Lambda),\quad
\mathcal D(\Omega)\widehat\otimes_i\mathcal D(\Lambda)
\simeq\mathcal D(\Omega\times\Lambda)
\]
łącznie z analogicznym izomorfizmem dla przestrzeni Schwartza,
porządkują związek między przestrzeniami testowymi a
ciągłymi formami dwuliniowymi. W szczególności nie wolno
bez dodatkowych założeń zastępować topologii $i$ przez $\pi$
w drugim wzorze.
\end{stwierdzenie}
\begin{proof}
Stwierdzenia o zanurzeniach i transpozycjach wynikają z
Twierdzenia~\ref{thm:zanurzenia25} i
Wniosku~\ref{cor:dual-lancuch25}.
Identyfikacje tensorowe udowodniono we wcześniejszym
rozdziale o produktach tensorowych; ich zestawienie
nie wymaga nowego przejścia granicznego.
\end{proof}

\begin{uwaga}[Co pozostaje do rozwinięcia]
\label{rem:future25}
Zastosowania do rozwiązywania równań różniczkowych
pozostają w odrębnym skrypcie. W tym opracowaniu
będziemy rozwijać dalej geometrię przestrzeni Banacha,
refleksywność i bazy Schaudera, własności aproksymacji,
teorię operatorów i analizę spektralną.
Moduły 15A (całkowanie w przestrzeniach lokalnie wypukłych)
i 15B (całki Bochnera, Pettisa, Dunforda i Gelfanda)
zostały już umieszczone bezpośrednio po \hyperref[sec:wyzsze-pochodne15]{rozdziale 15}.
Rozdział bieżący korzysta z wyników tych modułów jedynie
tam, gdzie jest to jawnie zaznaczone.
\end{uwaga}
